\documentclass[11pt]{article}
\usepackage[english]{babel}
\usepackage[a4paper,margin=0.9in]{geometry}
\usepackage[T1]{fontenc}
\usepackage[utf8]{inputenc}
\usepackage{lmodern}
\usepackage{amsmath,amssymb,mathtools}
\usepackage{microtype}
\usepackage{fancyhdr}
\usepackage{array,enumitem}
\setlength{\parindent}{1.4em}
\setlength{\parskip}{0.15em}
\emergencystretch=6em
\setlength{\headheight}{14pt}
\pagestyle{fancy}
\fancyhf{}
\cfoot{\thepage}
\newcommand{\sect}[2]{\section*{\S\,#1. #2}}
\lhead{Weber, Lehrbuch der Algebra I}
\rhead{English translation: \S\,62--\S\,68}
\begin{document}
\begin{center}
{\large Heinrich Weber}\\[0.35em]
{\Large Textbook of Algebra. Volume I.}\\[0.35em]
{\large Fifth Section. Linear Transformation.}\\[0.25em]
{\large \S\,62--\S\,68}
\end{center}

\sect{62}{Binary cubic forms}

A binary quadratic form gives no invariant constructions other than the discriminant. We shall therefore not concern ourselves with these, and pass at once to the consideration of the cubic forms
\begin{equation}
 f(x,y)=a_0x^3+a_1x^2y+a_2xy^2+a_3y^3.
\tag{1}
\end{equation}
Here we form, as the first quadratic covariant, the Hessian determinant. We shall always follow the rule here that the forms are written with integral numerical coefficients having no common factor. Thus in the determinant formed from the second derivatives of (1) we must discard the factor $4$, and obtain as first covariant
\[
 H=\begin{vmatrix}
 3a_0x+a_1y & a_1x+a_2y\\
 a_1x+a_2y & a_2x+3a_3y
 \end{vmatrix},
\]
or
\begin{equation}
 H=A_0x^2+A_1xy+A_2y^2,
\tag{2}
\end{equation}
where, for abbreviation,
\begin{equation}
\begin{aligned}
 A_0&=3a_0a_2-a_1^2,\\
 A_1&=9a_0a_3-a_1a_2,\\
 A_2&=3a_1a_3-a_2^2.
\end{aligned}
\tag{3}
\end{equation}

The discriminant of this quadratic form is an invariant of the given cubic form. It contains, however, the factor $3$ in all numerical coefficients, and therefore we put
\begin{equation}
 3D=4A_0A_2-A_1^2.
\tag{4}
\end{equation}
Then $D$ agrees with the discriminant of the cubic form [\S\,46, (10)]:
\begin{equation}
 D=a_1^2a_2^2+18a_0a_1a_2a_3-4a_0a_2^3-4a_1^3a_3-27a_0^2a_3^2.
\tag{5}
\end{equation}

We obtain a further cubic covariant by forming the functional determinant of $f(x,y)$ and $H(x,y)$:
\begin{equation}
 Q(x,y)=f'(x)H'(y)-f'(y)H'(x),
\tag{6}
\end{equation}
or
\begin{equation}
 Q=\begin{vmatrix}
 3a_0x^2+2a_1xy+a_2y^2 & 2A_0x+A_1y\\
 a_1x^2+2a_2xy+3a_3y^2 & A_1x+2A_2y
 \end{vmatrix},
\tag{7}
\end{equation}
from which no numerical factor can be removed. We set down here, in order, the computed coefficients of $x^3,x^2y,xy^2,y^3$, whose formation presents no difficulty:
\[
\begin{aligned}
 &27a_0^2a_3-9a_0a_1a_2+2a_1^3,\\
 &27a_0a_1a_3-18a_0a_2^2+3a_1^2a_2,\\
 &-27a_0a_2a_3+18a_1^2a_3-3a_1a_2^2,\\
 &-27a_0a_3^2+9a_1a_2a_3-2a_2^3.
\end{aligned}
\]

The behavior of $H,D,Q$ under a linear transformation follows from \S\,60, (2), (3). If $H',D',Q'$ are the corresponding constructions for a transformed form, then
\begin{equation}
 H'=r^2H,\qquad D'=r^6D,\qquad Q'=r^3Q.
\tag{8}
\end{equation}

We can use the covariants to transform the cubic form into a normal form which at the same time gives the solution of the cubic equation.

We choose the normal form
\begin{equation}
 f(x,y)=F(\xi,\eta)=\xi^3+\eta^3,
\tag{9}
\end{equation}
where $\xi,\eta$ are linear functions of $x,y$; as will appear at once, this form can always be produced when $D$ does not vanish, that is, when the equation $f=0$ does not have two equal roots. The solutions of $f=0$ then follow from the linear equations
\[
 \xi+\eta=0,
 \qquad \xi+\varepsilon\eta=0,
 \qquad \xi+\varepsilon^2\eta=0,
\]
where $\varepsilon$ is a third root of unity.

To form the forms $H',D',Q'$ for the normal form (9), we have to put $a'_0=a'_3=1$, $a'_1=a'_2=0$, and we obtain
\[
 H'=9\xi\eta,
 \qquad D'=-27,
 \qquad Q'=27(\xi^3-\eta^3).
\]
If the normal form (9) is derived from the general form $f(x,y)$ by a linear transformation, then by (8) one obtains
\begin{equation}
\begin{aligned}
 9\xi\eta&=r^2H,\\
 -27&=r^6D,\\
 27(\xi^3-\eta^3)&=r^3Q,\\
 \xi^3+\eta^3&=f.
\end{aligned}
\tag{10}
\end{equation}
It follows that
\[
 r^6=-\frac{27}{D},
 \qquad
 2\xi^3=\frac{r^3Q}{27}+f,
 \qquad
 2\eta^3=-\frac{r^3Q}{27}+f,
\]
and hence, after substituting for $r^3$ the value obtained from the first equation,
\begin{equation}
\begin{aligned}
 6\xi^3\sqrt{-3D}&=Q+3\sqrt{-3D}\,f,\\
 6\eta^3\sqrt{-3D}&=-Q+3\sqrt{-3D}\,f.
\end{aligned}
\tag{11}
\end{equation}
Thus $\xi$ and $\eta$ are determined by two cube roots, of which, by the first equation (10), one is determined by the other.

This also contains the proof that, when $D$ is different from zero, the normal form can be produced by linear transformation. Under this hypothesis $H$ splits into two distinct linear factors; and if, up to constant factors, these are chosen as the new variables $\xi,\eta$ of a linear transformation, then $A'_0=0$, $A'_2=0$. From the identical relations
\[
 3a'_3A'_0+a'_1A'_2=a'_2A'_1,
 \qquad
 a'_2A'_0+3a'_0A'_2=a'_1A'_1,
\]
it follows, unless $A'_1=0$ also, which is excluded by the non-vanishing of $D$, that $a'_1=0$, $a'_2=0$; that is, the transformed form of $f$ contains only the cubes of $\xi$ and $\eta$.

If one raises the first equation (10) to the third power and eliminates $\xi^3,\eta^3,r^6$ by means of (11) and of the second equation (10), one obtains the following identical relation among the covariants:
\begin{equation}
 4H^3+Q^2+27Df^2=0.
\tag{12}
\end{equation}

In order to carry out the transformation into the normal form, that is, actually to find the functions $\xi,\eta$, one decomposes the function $H$ into its linear factors:
\[
 4A_0H=(2A_0x+A_1y+\sqrt{-3D}\,y)(2A_0x+A_1y-\sqrt{-3D}\,y).
\]
Then $\xi,\eta$ differ from these two factors of $H$ only by one constant factor each. Thus, denoting two constant factors by $h,k$, we put
\begin{equation}
\begin{aligned}
 2\xi&=h(2A_0x+A_1y-\sqrt{-3D}\,y),\\
 2\eta&=k(2A_0x+A_1y+\sqrt{-3D}\,y),
\end{aligned}
\tag{13}
\end{equation}
from which, by multiplication and with reference to the first two equations (10), it follows that
\begin{equation}
 3hkA_0\sqrt[3]{-D}=1,
\tag{14}
\end{equation}
and the equations (11), by comparing the coefficients of $x^3$, give
\begin{equation}
\begin{aligned}
 6h^3A_0^3\sqrt{-3D}&=q_0+3\sqrt{-3D}\,a_0,\\
 6k^3A_0^3\sqrt{-3D}&=-q_0+3\sqrt{-3D}\,a_0,
\end{aligned}
\tag{15}
\end{equation}
where $q_0$ is the coefficient of $x^3$ in $Q$, namely
\begin{equation}
 q_0=27a_0^2a_3-9a_0a_1a_2+2a_1^3.
\tag{16}
\end{equation}

That the signs of $\sqrt{-3D}$ are to be taken as in formulas (13), with regard to the signs in (11), is seen most simply by carrying out the calculation for the special case $a_0=0$, $a_2=0$.

\sect{63}{The complete system of forms of the binary cubic form}

We shall still show that the invariants and covariants of the cubic form are exhausted by the functions $D,f,H,Q$, that is, that all invariants and covariants of a cubic form can be expressed rationally by these four forms.

Let
\[
 C=C(x,y,a)
\]
be a covariant of degree $\nu$ in the variables and of degree $\mu$ in the coefficients, and suppose that it has only integral numerical coefficients.

Then, for any transformed form,
\begin{equation}
 C'=C(\xi,\eta,a')=r^\lambda C(x,y,a),
\tag{1}
\end{equation}
where
\begin{equation}
 \lambda=\frac{3\mu-\nu}{2}.
\tag{2}
\end{equation}
We now understand by $\xi,\eta$ the variables of the normal form considered above. If $\varepsilon$ denotes a third root of unity, then
\[
 \xi=\varepsilon\xi',
 \qquad
 \eta=\varepsilon^2\eta'
\]
is a linear substitution with determinant $1$, by which the normal form \S\,62, (9) is carried into itself. Hence $C'$ cannot change when $\xi,\eta$ are replaced by $\varepsilon\xi,\varepsilon^2\eta$; therefore $C'$ can depend rationally only on $\xi\eta,\xi^3,\eta^3$.

The interchange of $\xi$ with $\eta$ corresponds to a substitution with determinant $-1$. Therefore $C'$ remains unchanged or changes its sign according as $\lambda$ is even or odd. Consequently $C'$ consists of a sum of terms of the form
\[
 M(\xi\eta)^\alpha(\xi^{3\beta}\pm\eta^{3\beta}),
\]
where $M$ is an integer factor, and the upper or lower sign applies according as $\lambda$ is even or odd. In the latter case this expression is divisible by $\xi^3-\eta^3$, and the quotient can be expressed rationally and integrally by $\xi^3+\eta^3$ and $\xi\eta$, since a symmetric function of two quantities can be represented in this way by their sum and product. Thus, putting $\rho=0$ or $=1$ according as $\lambda$ is even or odd, we obtain for $C'$ an expression of the form
\[
 C'(\xi,\eta)=(\xi^3-\eta^3)^\rho\sum M(\xi\eta)^\alpha(\xi^3+\eta^3)^\beta,
\]
where the $M$ are integral coefficients and $\alpha,\beta$ run through a series of positive integers satisfying
\begin{equation}
 3\rho+3\beta+2\alpha=\nu.
\tag{3}
\end{equation}
From (1), after multiplying by a sufficiently high power of $3$ and expressing $\xi\eta,\xi^3+\eta^3,\xi^3-\eta^3$ by $H,f,Q$ and $r$ by $D^{-1}$ according to \S\,62, (10), it follows that
\begin{equation}
 3^\sigma C=Q^\rho\sum MD^\gamma H^\alpha f^\beta,
\tag{4}
\end{equation}
where
\begin{equation}
 6\gamma=\lambda-3\rho-2\alpha,
\tag{5}
\end{equation}
and where the $M$ are again integral factors.

To every value of $\gamma$ there belongs, by (5), a definite value of $\alpha$, and by (3) a definite value of $\beta$; likewise, through a value of $\alpha$ or of $\beta$ the two other numbers are determined each time. Thus in (4) every exponent of $D,H$ or $f$ occurs only once.

That $\gamma$ is an integer follows easily from (2), (3), (5). Since by definition $\lambda-3\rho$ is an even number, $3\gamma$ is first of all an integer; furthermore, by (3),
\[
 \lambda-2\alpha=\frac32(\mu-\nu+2\beta+2\rho),
\]
and hence $6\gamma$, and therefore $3\gamma$, is divisible by $3$. It follows that $\gamma$ is also an integer.

That $\gamma$ must also be non-negative is seen as follows. Suppose that negative powers of $D$ occur on the right side of (4). Multiply the equation on both sides by so high a power of $D$ that no negative powers of $D$ appear on the right, but also so that the right side does not acquire the factor $D$. If in the equation so obtained we put
\begin{equation}
 a_0=0,
 \qquad a_1=1,
 \qquad a_2=0,
 \qquad a_3=0,
\tag{6}
\end{equation}
then
\begin{equation}
 D=0,
 \qquad H=-x^2,
 \qquad f=x^2y,
 \qquad Q=2x^3.
\tag{7}
\end{equation}
Thus the left side would vanish while the right side would not vanish, which is a contradiction.

Finally we can also prove that on the left side of (4) no power of $3$ can occur which is not contained in all factors $M$ and can therefore be cancelled.

In \S\,2 we proved the theorem that the product of two primitive integral rational functions of arbitrary variables is again a primitive function. By a primitive function there is meant one whose coefficients are integers without a common divisor. Here it remains to show that a sum of the form
\[
 Q^\rho\sum MD^\gamma H^\alpha f^\beta,
\]
where the $M$ are integers without a common divisor, cannot become imprimitive, and in particular cannot acquire the factor $3$, when the expressions in the $a,x,y$ are substituted for $Q,D,H,f$. By the theorem just mentioned, since $Q$ is a primitive function, it is enough to show this for the form
\begin{equation}
 \sum MD^\gamma H^\alpha f^\beta.
\tag{8}
\end{equation}
Here we may further omit all terms in which $M$ already has the factor $3$, and finally, again by the theorem just mentioned, we may suppose that the expression is not divisible by $D$, hence that it contains a term with $\gamma=0$. Assume then that, under these hypotheses, the developed expression (8) has the divisor $3$, and substitute the special values (6), (7), that is, put $D=0,H=-x^2,f=x^2y$. The expression then reduces to the single term in which $\gamma=0$, $2\alpha=\lambda$, namely
\[
 (-1)^\alpha M2^{\alpha+\rho}y^\beta,
\]
and it would follow that this $M$ must have the divisor $3$, contrary to the assumption. Thus we have proved:

Every integral covariant of the cubic form can be represented integrally and rationally, with integral coefficients, by one of the two expressions
\[
 \sum MD^\gamma H^\alpha f^\beta,
 \qquad
 Q\sum MD^\gamma H^\alpha f^\beta.
\]
The invariants are included among the covariants as a special case, and are all powers of $D$.

\sect{64}{Biquadratic forms}

We now pass to the consideration of the biquadratic form
\begin{equation}
 f(x,y)=a_0x^4+a_1x^3y+a_2x^2y^2+a_3xy^3+a_4y^4.
\tag{1}
\end{equation}

We first have the Hessian determinant as a covariant of the fourth order, which, in order to free it from a numerical factor, we divide by $3$:
\begin{equation}
 H=\frac13
 \begin{vmatrix}
 12a_0x^2+6a_1xy+2a_2y^2 & 3a_1x^2+4a_2xy+3a_3y^2\\
 3a_1x^2+4a_2xy+3a_3y^2 & 2a_2x^2+6a_3xy+12a_4y^2
 \end{vmatrix},
\tag{2}
\end{equation}
or
\[
 H=A_0x^4+A_1x^3y+A_2x^2y^2+A_3xy^3+A_4y^4,
\]
where
\[
\begin{aligned}
 A_0&=8a_0a_2-3a_1^2,& A_4&=8a_2a_4-3a_3^2,\\
 A_1&=24a_0a_3-4a_1a_2,& A_3&=24a_1a_4-4a_2a_3,\\
 A_2&=48a_0a_4+6a_1a_3-4a_2^2,
\end{aligned}
\]
and a covariant of the sixth degree
\begin{equation}
 T=\frac{1}{12}
 \begin{vmatrix}
 f'(x)&f'(y)\\
 H'(x)&H'(y)
 \end{vmatrix}.
\tag{3}
\end{equation}

The invariants of the biquadratic form are most simply formed from the root differences, according to \S\,61.

We thus obtain an invariant of the second and one of the third order in the coefficients:
\begin{equation}
 A=\frac12 a_0^2\bigl[(12)^2(34)^2+(13)^2(24)^2+(14)^2(23)^2\bigr],
\tag{4}
\end{equation}
\begin{equation}
 B=a_0^3\sum(12)^2(34)^2(13)(42).
\tag{5}
\end{equation}
These expressions may be written more clearly if one puts
\begin{equation}
 U=(12)(34),\qquad V=(13)(42),\qquad W=(14)(23),
\tag{6}
\end{equation}
namely
\begin{equation}
\begin{aligned}
 A&=\frac12 a_0^2(U^2+V^2+W^2),\\
 B&=a_0^3\{U^2(V-W)+V^2(W-U)+W^2(U-V)\}\\
  &=a_0^3(W-V)(U-W)(V-U).
\end{aligned}
\tag{7}
\end{equation}

The necessary formulas for computing these quantities were already developed in \S\,47. From the formulas (10) there one obtains
\[
 U=v^2-w^2,
 \qquad V=w^2-u^2,
 \qquad W=u^2-v^2,
\]
where $u^2,v^2,w^2$ are the roots of the cubic resolvent
\[
 z^3+2az^2+(a^2-4c)z-b^2=0.
\]
It follows that
\[
\begin{aligned}
 A&=a_0^2\{(u^2+v^2+w^2)^2-3(v^2w^2+w^2u^2+u^2v^2)\}\\
  &=-72ac+27b^2,\\[0.4em]
 B&=a_0^3(2u^2-v^2-w^2)(2v^2-w^2-u^2)(2w^2-u^2-v^2)\\
  &=a_0^3(3u^2+2a)(3v^2+2a)(3w^2+2a)\\
  &=a_0^3(2a^3-72ac+27b^2),
\end{aligned}
\]
so that $A,B$ have the same meaning as in \S\,47, and, expressed by the coefficients $a_0,a_1,a_2,a_3,a_4$, are represented as follows:
\begin{equation}
 A=a_2^2-3a_1a_3+12a_0a_4,
\tag{8}
\end{equation}
\begin{equation}
 B=27a_0a_3^2+27a_1^2a_4+2a_2^3-72a_0a_2a_4-9a_1a_2a_3.
\tag{9}
\end{equation}
We had also already formed there the discriminant of the biquadratic equation as a third invariant, but one depending on $A$ and $B$,
\begin{equation}
 D=a_0^6U^2V^2W^2,
\tag{10}
\end{equation}
and found
\begin{equation}
 27D=4A^3-B^2.
\tag{11}
\end{equation}

From the formula of \S\,67, $2\lambda=n\mu-\nu$, we obtain for the invariant constructions found so far, which we denote by accents for a transformed form, the following relations:
\begin{equation}
 H'=r^2H,
 \qquad T'=r^3T,
\tag{12a}
\end{equation}
\begin{equation}
 A'=r^4A,
 \qquad B'=r^6B,
 \qquad D'=r^{12}D.
\tag{12b}
\end{equation}

\sect{65}{Solution of the biquadratic equation}

We want to produce a normal form by a linear substitution with determinant $1$. For its derivation we assume the given biquadratic form to be decomposed into linear factors:
\begin{equation}
 f(x,y)=a_0(x-\alpha y)(x-\beta y)(x-\gamma y)(x-\delta y).
\tag{1}
\end{equation}
We put
\begin{equation}
\begin{aligned}
 x\sqrt{\alpha-\beta}&=\beta\xi-\alpha\eta,\\
 y\sqrt{\alpha-\beta}&=\xi-\eta,
 \qquad r=1,
\end{aligned}
\tag{2}
\end{equation}
so that
\begin{equation}
\begin{aligned}
 x-\alpha y&=-\sqrt{\alpha-\beta}\,\xi,\\
 x-\beta y&=-\sqrt{\alpha-\beta}\,\eta,\\
 (x-\gamma y)\sqrt{\alpha-\beta}&=(\beta-\gamma)\xi-(\alpha-\gamma)\eta,\\
 (x-\delta y)\sqrt{\alpha-\beta}&=(\beta-\delta)\xi-(\alpha-\delta)\eta.
\end{aligned}
\tag{3}
\end{equation}
Thus $f(x,y)$ goes over into
\begin{equation}
 F(\xi,\eta)=a'_1\xi^3\eta+a'_2\xi^2\eta^2+a'_3\xi\eta^3.
\tag{4}
\end{equation}
By (3), the coefficients $a'_1,a'_2,a'_3$ are easily expressed through $\alpha,\beta,\gamma,\delta$:
\begin{equation}
\begin{aligned}
 a'_1&=a_0(\beta-\gamma)(\beta-\delta),\\
 a'_3&=a_0(\alpha-\gamma)(\alpha-\delta),\\
 a'_2&=-a_0\{(\alpha-\gamma)(\beta-\delta)+(\alpha-\delta)(\beta-\gamma)\}.
\end{aligned}
\tag{5}
\end{equation}
By interchanging the roots we can produce the normal form (4) in six different ways, but these yield only three distinct values of $a'_2$. These three values are, if as above
\begin{equation}
\begin{aligned}
 U&=(\alpha-\beta)(\gamma-\delta),\\
 V&=(\alpha-\gamma)(\delta-\beta),\\
 W&=(\alpha-\delta)(\beta-\gamma),\\
 U+V+W&=0,
\end{aligned}
\tag{6}
\end{equation}
then
\begin{equation}
 a'_2=a_0(V-W),
 \qquad a''_2=a_0(W-U),
 \qquad a'''_2=a_0(U-V).
\tag{7}
\end{equation}
If $a'_2,a''_2,a'''_2$ are known, then $U,V,W$ can also be determined by
\begin{equation}
\begin{aligned}
 3a_0U&=a'''_2-a''_2,\\
 3a_0V&=a'_2-a'''_2,\\
 3a_0W&=a''_2-a'_2.
\end{aligned}
\tag{8}
\end{equation}
It follows that, if of the three quantities $a'_2,a''_2,a'''_2$ two are equal, one of the quantities $U,V,W$ vanishes, so that two of the roots of the biquadratic equation are equal, and consequently its discriminant vanishes.

Let us also put
\begin{equation}
\begin{aligned}
 \alpha\beta+\gamma\delta&=u,\\
 \alpha\gamma+\delta\beta&=v,\\
 \alpha\delta+\beta\gamma&=w.
\end{aligned}
\tag{9}
\end{equation}
Then
\[
 a_0(u+v+w)=a_2
\]
[the sum of the products of the roots of $f(x,y)$ taken two at a time], and
\[
 U=v-w,
 \qquad V=w-u,
 \qquad W=u-v.
\]
Thus
\begin{equation}
\begin{aligned}
 3a_0u&=a_2-a'_2,\\
 3a_0v&=a_2-a''_2,\\
 3a_0w&=a_2-a'''_2.
\end{aligned}
\tag{10}
\end{equation}
Consequently, together with the quantities $a'_2,a''_2,a'''_2$, the quantities $u,v,w$ are also known. But if one knows $u,v,w$, then the solution of the biquadratic equation has been reduced to square roots. For
\[
 \alpha\beta+\gamma\delta=u,
 \qquad a_0\alpha\beta\gamma\delta=a_4,
\]
and therefore
\[
 2\alpha\beta=u+\sqrt{\frac{a_0u^2-4a_4}{a_0}},
 \qquad
 2\gamma\delta=u-\sqrt{\frac{a_0u^2-4a_4}{a_0}}.
\]
Since the products
\[
 \alpha\gamma,
 \qquad \alpha\delta,
 \qquad \beta\gamma,
 \qquad \beta\delta
\]
can likewise be determined, $\alpha^2$, for example, can be computed from
\[
 \frac{\alpha\beta\cdot\alpha\gamma}{\beta\gamma}.
\]

It remains only to form the cubic equation whose roots are $a'_2,a''_2,a'''_2$. This is obtained at once from the invariants.

If the invariants $A',B'$ are formed from the transformed form (4), then by \S\,64, (8), (9), (12), since $r=1$, one obtains
\begin{equation}
\begin{aligned}
 A&=a_2'^2-3a'_1a'_3,\\
 B&=2a_2'^3-9a'_1a'_2a'_3,\\
 D&=a_1'^2a_2'^2a_3'^2-4a_1'^3a_3'^3.
\end{aligned}
\tag{11}
\end{equation}
Eliminating $a'_1a'_3$ gives
\[
 a_2'^3-3a'_2A+B=0.
\]
Hence $a'_2,a''_2,a'''_2$ are the roots of the cubic equation in $z$
\begin{equation}
 z^3-3Az+B=0.
\tag{12}
\end{equation}
This is probably the simplest form that can be given to the cubic resolvent of the biquadratic equation.

\sect{66}{The covariants}

If we form the covariant $H$, according to \S\,64, (2), for the transformed form
\begin{equation}
 f(x,y)=F(\xi,\eta)=a'_1\xi^3\eta+a'_2\xi^2\eta^2+a'_3\xi\eta^3,
\tag{1}
\end{equation}
we obtain
\begin{equation}
 -H=3a_1'^2\xi^4+4a'_1a'_2\xi^3\eta-(6a'_1a'_3-4a_2'^2)\xi^2\eta^2+4a'_2a'_3\xi\eta^3+3a_3'^2\eta^4,
\tag{2}
\end{equation}
and from this it follows that
\begin{equation}
 H+4a'_2f=-3(a'_1\xi^2-a'_3\eta^2)^2.
\tag{3}
\end{equation}
Thus, apart from the factor $-3$, this combination is the square of a quadratic form; the same holds when $a'_2$ is replaced by $a''_2,a'''_2$. For abbreviation, we therefore put
\begin{equation}
\begin{aligned}
 H+4a'_2f&=-3\psi_1^2,\\
 H+4a''_2f&=-3\psi_2^2,\\
 H+4a'''_2f&=-3\psi_3^2.
\end{aligned}
\tag{4}
\end{equation}

Of the three functions $\psi_1,\psi_2,\psi_3$, no two can have a common divisor if the discriminant $D$ is assumed different from zero. For then the three quantities $a'_2,a''_2,a'''_2$ are also distinct. Thus, if $\psi_1$ and $\psi_2$ vanish, then $H$ and $f$ must vanish; that is, a common divisor of $\psi_1$ and $\psi_2$ would have to be a common divisor of $H$ and $f$. But $\xi$ was an arbitrary linear divisor of $f$. By (2), this can be a divisor of $H$ only when $a'_3=0$, hence again only when the discriminant $D$ vanishes [\S\,65, (11)].

Now the functional determinant $T$ of $f$ and $H$ is identical with the functional determinant of $f$ and $H+\lambda f$, whatever $\lambda$ may be; and hence, putting $\lambda=4a'_2$, one obtains
\[
 T=-\frac13\psi_1\bigl[F'(\xi)\psi_1'(\eta)-F'(\eta)\psi_1'(\xi)\bigr].
\]
Thus $T$ is divisible by $\psi_1$, and by the same reasoning by $\psi_2$ and $\psi_3$; consequently $T$ is, up to a factor independent of $\xi,\eta$, identical with the product $\psi_1\psi_2\psi_3$.

If we now form the product of the three equations (4), then, with reference to the cubic equation \S\,65, (12), satisfied by $a'_2,a''_2,a'''_2$, it follows that
\[
 H^3-48AHf^2-64Bf^3=cT^2,
\]
where $c$ is a constant still to be determined. By \S\,64, (12), $c$ remains unchanged under a linear transformation, and therefore its value is found by equating, in the normal form (1), (2), the terms with the highest power of $\xi$ on both sides: $c=-27$.

We thus have the following identical relation among the invariants and covariants:
\begin{equation}
 H^3-48AHf^2-64Bf^3=-27T^2.
\tag{5}
\end{equation}

The functions $\psi_1,\psi_2,\psi_3$ are likewise covariants, although with coefficients not rational in the coefficients of $f$. We can express them easily in terms of the roots $\alpha,\beta,\gamma,\delta$ of $f=0$. Putting $y=1$ for simplicity, it follows from \S\,65, (3) and (5) that
\[
 \psi_1=a'_1\xi^2-a'_3\eta^2
 =a_0\frac{(\beta-\gamma)(\beta-\delta)(x-\alpha)^2-(\alpha-\gamma)(\alpha-\delta)(x-\beta)^2}{\alpha-\beta}.
\]
Here, however, $\alpha-\beta$ cancels from numerator and denominator, and $\psi_1$ can be represented in two forms; to obtain $\psi_2$ and $\psi_3$, one then only has to cyclically permute $\beta,\gamma,\delta$.

One finds
\begin{equation}
\begin{aligned}
 \frac{\psi_1}{a_0}
 &= (\gamma-\beta)(x-\alpha)(x-\delta)-(\alpha-\delta)(x-\beta)(x-\gamma)\\
 &= (\delta-\beta)(x-\alpha)(x-\gamma)-(\alpha-\gamma)(x-\beta)(x-\delta),\\[0.4em]
 \frac{\psi_2}{a_0}
 &= (\delta-\gamma)(x-\alpha)(x-\beta)-(\alpha-\beta)(x-\gamma)(x-\delta)\\
 &= (\beta-\gamma)(x-\alpha)(x-\delta)-(\alpha-\delta)(x-\gamma)(x-\beta),\\[0.4em]
 \frac{\psi_3}{a_0}
 &= (\beta-\delta)(x-\alpha)(x-\gamma)-(\alpha-\gamma)(x-\delta)(x-\beta)\\
 &= (\gamma-\delta)(x-\alpha)(x-\beta)-(\alpha-\beta)(x-\delta)(x-\gamma).
\end{aligned}
\tag{6}
\end{equation}

The quantities $\psi_1^2,\psi_2^2,\psi_3^2$ may be regarded as the roots of a cubic equation whose coefficients are covariants. From (4) this equation is readily obtained in the form
\[
 \left(z+\frac{H+4a'_2f}{3}\right)
 \left(z+\frac{H+4a''_2f}{3}\right)
 \left(z+\frac{H+4a'''_2f}{3}\right)=0.
\]
By \S\,65, (12), and using expression (5) for $T^2$, this equation gives
\begin{equation}
 z^3+Hz^2+\frac13(H^2-16Af^2)z-T^2=0.
\tag{7}
\end{equation}

It is still of interest to form the discriminant of this equation. This is obtained most simply from the expression
\[
 \Delta=(\psi_1^2-\psi_2^2)^2(\psi_1^2-\psi_3^2)^2(\psi_2^2-\psi_3^2)^2,
\]
when one substitutes the expressions (4) and then replaces $(a'_2-a''_2)^2(a'_2-a'''_2)^2(a''_2-a'''_2)^2$ by the discriminant of the equation \S\,65, (12), where $D$ then denotes the discriminant of $f$. One obtains
\begin{equation}
 \Delta=2^{12}Df^6.
\tag{8}
\end{equation}

\sect{67}{The complete invariant system of the binary biquadratic form}

We shall still prove that the constructions $A,B$ exhaust the system of independent invariants of the binary biquadratic form. While one is usually content to show that every invariant is an integral rational function of $A,B$, we shall, as we have already done in \S\,63 for the covariants of the cubic form, also touch upon the question of the numerical coefficients; here a new phenomenon appears.

For example, by \S\,64, (11), the invariant $D$ is indeed rationally expressible through $A,B$. But the numerical coefficients are not integers; rather, they have denominator $27$, although the coefficients in the developed function $D$ are all integers.

We therefore now consider as integral invariants those integral rational homogeneous functions of the $\mu$th degree in the five variables $a_0,a_1,a_2,a_3,a_4$
\begin{equation}
 I(a_0,a_1,a_2,a_3,a_4)=I(a),
\tag{1}
\end{equation}
whose coefficients are integers, and which have the invariant property
\begin{equation}
 I'=r^{2\mu}I,
\tag{2}
\end{equation}
where $I'$ or $I(a')$ is the same function of the coefficients of a transformed function.

Such invariants are $A,B,D$, between which the relation
\begin{equation}
 27D=4A^3-B^2
\tag{3}
\end{equation}
holds; and our goal is to prove that all integral invariants are integral rational functions of these three, with integer coefficients.

If in the normal form (4) of \S\,65 we form the function $I'$, then from (2), since $r=1$, we first obtain
\[
 I(a)=I'(0,a'_1,a'_2,a'_3,0).
\]
This function $I$ can depend only on $a'_2$ and on the product $a'_1a'_3$, for the substitution
\[
 \xi=\lambda\xi',
 \qquad
 \eta=\frac1\lambda\eta'
\]
with determinant $1$ leaves the normal form \S\,65, (4), and in it $a'_2$, unchanged, while $a'_1,a'_3$ go over into $\lambda^2a'_1$, $a'_3:\lambda^2$; here $\lambda$ is arbitrary.

Thus, putting
\[
 a'_2=z,
\]
we have
\begin{equation}
 I=\varphi(a'_1a'_3,z),
\tag{4}
\end{equation}
where $\varphi$ is an integral rational function with integer coefficients of the two arguments $a'_1a'_3$ and $z$.

But by \S\,65, (11),
\[
 3a'_1a'_3=z^2-A.
\]
Therefore, if we multiply (4) by a suitable power of $3$, say $3^\nu$, we obtain
\begin{equation}
 3^\nu I=\psi(A,z),
\tag{5}
\end{equation}
where $\psi$ is again an integral function of $A$ and $z$. But by \S\,65, (12),
\[
 z^3=3Az-B,
\]
and hence all higher powers of $z$ can be expressed by the first and second powers. Thus we obtain
\begin{equation}
 3^\nu I=\chi(A,B,z),
\tag{6}
\end{equation}
where $\chi$ is of degree at most two with respect to $z$.

The left side of (6), however, remains unchanged when $z$ is replaced by $a'_2,a''_2,a'''_2$, that is, by three different values. Consequently the function $\chi$ cannot contain the variable $z$ at all (\S\,30, II).

Thus
\begin{equation}
 3^\nu I=\chi(A,B),
\tag{7}
\end{equation}
where $\chi$ is an integral, rational function with integer coefficients.

If in (7) we put, according to (3),
\[
 B^2=4A^3-27D,
\]
it follows that (7) has one of the two following forms:
\begin{equation}
 3^\nu I=\Phi(A,D),
 \qquad
 B\Phi(A,D),
\tag{8}
\end{equation}
according as the degree of $I$ is even or odd.

If $I$ is a primitive function of the $a$, that is, one whose numerical coefficients have no common divisor, then the coefficients of the functions $\Phi$ can in any case have no common divisor other than a power of $3$.

What remains to be proved is that they are all divisible by $3^\nu$, or, what is the same, that if $\Phi(A,D)$ is assumed primitive, then $\nu=0$.

The question is therefore this: can functions with divisor $3$ arise from the primitive function $\Phi(A,D)$ or $B\Phi(A,D)$ when $A,B,D$ are replaced by their expressions in the $a$? Since $B$ is primitive, by \S\,2 the function $B\Phi(A,D)$ expressed in the $a$ can have the divisor $3$ only if $\Phi(A,D)$ has it.

Thus the question is reduced to this: can the primitive function $\Phi(A,D)$ acquire the divisor $3$ when $A,D$ are replaced by their expressions?

That this question must be answered in the negative can be seen easily. First imagine that in $\Phi(A,D)$ all terms whose coefficients are divisible by $3$ have been removed; for plainly the remaining terms must still have the same property, namely, to acquire the divisor $3$ by substitution of the expressions for $A,D$.

Second, we may assume that $\Phi(A,D)$ does not have the factor $D$, for by omitting this factor, according to the theorem already mentioned (\S\,2), the property in question would not be destroyed.

Then, however, a number divisible by $3$ would have to arise from $\Phi(A,D)$ when all $a_i$, except $a_2$, are set equal to zero and $a_2=1$ is put. This gives $D=0$, $A=1$, and therefore $\Phi(1,0)$ would have to be a number divisible by $3$. But this is the coefficient of the highest power of $A$ in $\Phi(A,D)$, which by assumption is not divisible by $3$. Thus we have proved:

Every integral invariant of the biquadratic form can be represented rationally and integrally by $A,B,D$, and indeed in one of the two forms
\[
 \Phi(A,D),
 \qquad
 B\Phi(A,D).
\]
Conversely, it is clear by itself that every such expression, if homogeneous in the coefficients $a$, represents an integral invariant\footnote{The theory of invariants is presented in detail in the works: Clebsch, \emph{Theory of Binary Algebraic Forms}. Leipzig 1872. Faa di Bruno, \emph{Introduction to the Theory of Binary Forms}. German translation by Walter. Leipzig 1881. P. Gordan, \emph{Lectures on Invariant Theory}, edited by Kerschensteiner, Leipzig 1885. See also Franz Meyer, \emph{Report on the Present State of Invariant Theory}, in the Annual Report of the German Mathematical Society 1890/91 (Berlin 1892).}.

\begin{center}
{\large Sixth Section.}\\[0.45em]
{\Large Tschirnhausen Transformation.}
\end{center}

\sect{68}{Hermite's form of the Tschirnhausen transformation}

In the fourth section we have already become acquainted with the basic idea of the Tschirnhausen transformation.

The problem was to transform an algebraic equation of degree $n$,
\begin{equation}
 f(x)=a_0x^n+a_1x^{n-1}+\cdots+a_{n-1}x+a_n=0,
\tag{1}
\end{equation}
by a substitution of degree $n-1$,
\begin{equation}
 y=\alpha_0+\alpha_1x+\alpha_2x^2+\cdots+\alpha_{n-1}x^{n-1},
\tag{2}
\end{equation}
so as to obtain, in the arbitrary coefficients $\alpha$, means for simplifying the equation.

Hermite greatly simplified this problem and brought it into connection with invariant theory by giving the substitution (2) a special form\footnote{Hermite, ``Sur quelques théorèmes d'algèbre et la résolution de l'équation du quatrième degré,'' separately issued from the Comptes rendus of the Paris Academy. Paris 1859.}.

In \S\,4 we became acquainted with a sequence of functions $f_0,f_1,f_2,\ldots,f_{n-1}$, by means of which the powers of $x$ up to the $(n-1)$st can be expressed rationally; these are the functions which Hermite uses to represent the substitution (2).

Here $x$ denotes a root of the equation (1), while $t$ is to be understood as an indeterminate variable. Then by \S\,4,
\begin{equation}
 \frac{f(t)}{t-x}=t^{n-1}f_0(x)+t^{n-2}f_1(x)+\cdots+t f_{n-2}(x)+f_{n-1}(x)
\tag{3}
\end{equation}
and
\begin{equation}
\begin{aligned}
 f_0(x)&=a_0,\\
 f_1(x)&=a_0x+a_1,\\
 f_2(x)&=a_0x^2+a_1x+a_2,\\
 &\hspace{2.5em}\cdots\\
 f_{n-1}(x)&=a_0x^{n-1}+a_1x^{n-2}+a_2x^{n-3}+\cdots+a_{n-1}.
\end{aligned}
\tag{4}
\end{equation}
We now take the substitution (2) in the form
\begin{equation}
 y=t_{n-1}f_0(x)+t_{n-2}f_1(x)+\cdots+t_1f_{n-2}(x)+t_0f_{n-1}(x),
\tag{5}
\end{equation}
where $t_{n-1},t_{n-2},\ldots,t_1,t_0$ are the indeterminate quantities which have taken the place of the $\alpha$, and are not to be confused with the powers of $t$. But $y$ is obtained from the expression (3) when $t^k$ is replaced by $t_k$.

As in \S\,42, let the sign $S$ placed before a function mean that the sum is to be taken over all roots $x$ of equation (1). Then
\begin{equation}
\begin{aligned}
 S[f_0(x)]&=na_0,\\
 S[f_1(x)]&=(n-1)a_1,\\
 S[f_2(x)]&=(n-2)a_2,\\
 &\hspace{2.5em}\cdots\\
 S[f_{n-1}(x)]&=a_{n-1},
\end{aligned}
\tag{6}
\end{equation}
and hence
\begin{equation}
 S(y)=na_0t_{n-1}+(n-1)a_1t_{n-2}+\cdots+2a_{n-2}t_1+a_{n-1}t_0,
\tag{7}
\end{equation}
an expression which is obtained from $f'(t)$ by replacing $t^k$ by $t_k$.

If with the help of (7) we eliminate the variable $t_{n-1}$ from (5), we get
\begin{equation}
 y-\frac1nS(y)=t_{n-2}F_0(x)+t_{n-3}F_1(x)+\cdots+t_0F_{n-2}(x),
\tag{8}
\end{equation}
where
\begin{equation}
\begin{aligned}
 F_0&=f_1-\frac{n-1}{n}a_1=a_0x+\frac{a_1}{n},\\
 F_1&=f_2-\frac{n-2}{n}a_2=a_0x^2+a_1x+\frac{2a_2}{n},\\
 &\hspace{2.5em}\cdots\\
 F_{n-2}&=f_{n-1}-\frac1n a_{n-1}=a_0x^{n-1}+a_1x^{n-2}+\cdots+\frac{n-1}{n}a_{n-1}.
\end{aligned}
\tag{9}
\end{equation}
Thus
\begin{equation}
 S[F_0(x)]=0,
 \quad S[F_1(x)]=0,
 \ldots,
 \quad S[F_{n-2}(x)]=0.
\tag{10}
\end{equation}
Putting, according to (3),
\begin{equation}
\begin{aligned}
 F(t,x)&=\frac{f(t)}{t-x}-\frac1n f'(t)\\
 &=t^{n-2}F_0(x)+t^{n-3}F_1(x)+\cdots+tF_{n-3}(x)+F_{n-2}(x),
\end{aligned}
\tag{11}
\end{equation}
the right side of (8) is obtained from $F(t,x)$ by replacing $t^k$ by $t_k$.

If from the outset we take $y$ in the form
\begin{equation}
 y=t_{n-2}F_0(x)+t_{n-3}F_1(x)+\cdots+t_0F_{n-2}(x),
\tag{12}
\end{equation}
then the equation
\begin{equation}
 S(y)=0
\tag{13}
\end{equation}
is identically satisfied. What use this form of the substitution affords will be shown by the next considerations.

\end{document}
